\usection{Virtù Creative}

Ho menzionato in precedenza il quadro tecnico generale del gioco in questione, ma proviamo ad andare più a fondo. Le regole e le procedure sono, alla fin fine, vuote e fallibili a meno che non siano guidate dalla comprensione dell'io creativo. Immagina di partecipare a una partita di football senza cercare di vincere: giochi seguendo tutte le regole e nessuno può indicare una cosa specifica che stai facendo male, perché l'errore che rende il gioco un rituale vuoto è nello spirito dell'esercizio stesso. Questa è la differenza che fanno le virtù creative.

Come discusso nella sezione base del libro, l'ethos creativo del \wargam{ing} riguarda la sportività baldanzosa e la curiosità erudita: prepariamo scenari e vi scateniamo contro i giocatori per imparare e crescere. Le virtù creative sono l'elemento distinto che orienta l'attività in modo da favorire questi obiettivi creativi.

Simili virtù creative si sviluppano solo col tempo e la pratica. Non è inusuale per una partita, o un individuo, iniziare senza simili astruse considerazioni. Sul lungo termine, però, ciò che ci fa continuare a giocare una volta che si perde l'effetto novità è la sostanza creativa coltivata dal gioco virtuoso. Non mi aspetto che le mie divagazioni sulle virtù creative abbiano molto senso per tutti i lettori, ma posso speare di fornire qualche nome per delle idee che anche altri hanno percepito nel gioco.

\usubsection{L'autenticità ha valore}

L'hobby lotta da tempo con i desideri contrastanti di essere, da una parte, la pratica di un'abilità e, dall'altra, un processo di consumo di contenuto. Nel profondo della sua essenza, il gioco è un'esperienza che compri e esperisci, o è un'abilità che pratichi e in cui migliori? L'industria del gioco è ampiamente focalizzata sul primo caso: la maggior parte dei prodotti di gioco non cercano di dirti come giocare, piuttosto si offrono di essere un contenuto che consumi, in un modo o nell'altro.

Credo che lo scopo di giocare secondo la via del \wargam{ing} sia di mettere in piedi uno scenario significativo e risolverlo usando la saggezza e il giudizio del gruppo. Quando succede, lo chiamo "autenticità". All'opposto del gioco autentico c'è l'esperire quella che alla fin fine è una presentazione mediatica preparata in anticipo. Può essere il Game Master che la prepara in anticipo, o può persino essere il produttore del gioco che vende materiale che verrà poi presentato dal GM. In entrambi i casi, l'attività non è autentica e non coinvolge vero pensiero, vere emozioni, vera elaborazione e un effettivo "giocarsela" tra scenario e giocatori.

Essere consci dell'autenticità significa un cambiamento radicale nel modo in cui si pensa all'intero edificio del gioco, perché molti metodi e schemi di pensiero comunamente suggeriti sono profondamente non-autentici. Imparare a prestare attenzione nel momento in cui i dadi colpiscono il tavolo a se i giocatori stiano \textbf{esperendo} e \textbf{partecipando} aiuta nella ricerca dell'autenticità.

È anti-autentico che l'arbitro si preoccupi di far seguire al gioco sentieri predefiniti diretti a risultati predeterminati. Questo si chiama spesso "railroading", ma non credo che sia utile etichettare, demonizzare e mettere da parte qualcosa che è un conflitto di valori. Le pratiche dell'arbitro possono facilmente diventare non autentiche molto prima che comincino ad assomigliare al \textit{railroading} tradizionale. Dargli un nome è, talvolta, un modo di fingere che tu non lo stia facendo.

Anche i giocatori possono essere più o meno autentici. È autentico prendere per buona l'autenticità del gioco, ingaggiandolo allo stesso livello; è non autentico essere ironici o cinici, credendo che il gioco sia in realtà truccato, in verità un'opera teatrale manipolata da un arbitro che si suppone imparziale. È autentico fare del tuo meglio; è non autentico buttare via la partita. È autentico prendersi la responsabilità per il destino del tuo personaggio e del suo gruppo di avventurieri; è non autentico credere che sarà il processo del gioco a trasportatvi e che le vostre scelte non determinano il vostro destino. È autentico ridere e piangere in risposta al gioco.

\dnd{} è un \textit{franchise} culturale che è diventato spaventosamente non autentico nel processo che l'ha avvolto in un comodo prodotto di consumo. Scegliete il contenuto della vostra partita, e lo risolvete da voi, o ci sono forze che fanno queste cose per voi? Vi importa?

\usubsection{Arbitraggio egalitario}

L'arbitraggio di \dnd{} ha, nel suo retaggio culturale, tratti sociologici da imbonitore. C'è, sottostante, l'idea che l'arbitro sia un incredibile mago che \textit{lascia} che i giocatori entrino nel suo mondo magico. Le sue vie sono misteriose, le regole del gioco cangianti e i giudizi arcani. Gli altri giocatori sono incoraggiati a trattare l'arbitro come una figura autoritaria.

A posteriori, mi sembra che il modello dell'arbitro-imbonitore fosse e sia percepito come utile in parte per gli aspetti più legati alle capacità di intrattenitore e in parte perché permette all'arbitro di nascondere la preparazione. Il gioco può essere molto più impressionante con meno lavoro se l'arbitro viene messo su un piedistallo come questo. Ha anche aiutato a stabilire un'\textbf{organizzazione autoritaria} del gruppo di gioco, che è meglio di nessuna organizzazione.

Ma, d'altra parte, ha anche dei problemi. Le dinamiche da ciarlatano incoraggiano l'infantilizzazione dei giocatori, l'imbroglio dell'arbitro e la mancanza di responsabilità a lungo termine. E non dimentichiamo le lotte di potere, l'idea di avversarietà che può venir associata alla posizione di arbitro.

Per queste ragioni, ho concluso che non posso portare il gioco dove voglio essendo il misterioso Signore del Gioco\translatornote{In inglese era "Master of the Game", che riprende il termine Game Master usato spesso per indicare l'arbitro.}. Devo essere meno, cosicché i giocatori possano crescere.

L'alternativo \textbf{ideale egalitario} è di trattare una campagna di \wargam{ing} come un progetto collettivo intrapreso da tutti i partecipanti come socialmente pari. C'è un arbitro-GM, che è molto probabilmente il motore primo creativo del progetto, ma non desidera che gli altri giocatori entrino nel suo mondo \textit{per avere un pubblico}: l'arbitro di \wargam{e} desidera partecipare in un \textbf{discorso tra pari} che non può avvenire fintantoché tu sei il mago e gli altri dei cafoni.

Per essere un arbitro egalitario, credo sia importante praticare l'\textbf{emancipazione del giocatore}, che significa aiutare gli altri partecipanti ad avviarsi da sé come giocatori: insegnagli le regole, insegnagli la teoria dell'arte, incoraggia il percorso creativo e la passione di ciascuno dei tuoi co-giocatori. Finché sei tu il più grande, questa è una tua responsabilità sociale indisputabile in quanto arbitro; se vuoi liberarti di questo fardello, concludi il processo e siate uguali.

Devono esserci \textbf{trasparenza} e \textbf{responsabilità}, come in qualsiasi organizzazione costituita. L'arbitro è volenteroso e in grado di giustificare le proprie decisioni e di sviluppare la campagna nelle direzioni desiderate dai suoi soci nel gioco. Attitudini simili aiutano a far crescere la fiducia e, con la fiducia, la vede nel gioco e l'eccitazione nei risultati della campagna. Questo è un antidoto diretto al cinismo e incoraggia una partecipazione autentica al gioco.

Partecipare seriamente dei desiderei creativi e dell'ideazione di altri partecipanti significa spesso fare le cose "male": dei due, sei spesso il più esperto ed educato nell'arte, quindi è naturale. Comunque, non penso che sia una buona idea trarne come conseguenza che l'arbitro deve sempre guidare e gli altri seguire. La differenza tra gioco e lezion e come metodi di apprendimento è che il gioco è un'attività. Si impara meglio se lasci spazio nella campagna per provare idee e approcci. Discutetele, certo; spiega cosa ne pensi e perché vorresti fare le cose in un modo o un altro; alla fine, però, prova il loro modo con cuore onesto, lascia che vediate tutti come va e, forse, crescerai tu stesso nello scambio.

\usubsection{Uno scambio giudiziale negoziato}

L'applicazione della filosofia egalitaria ai dettagli dell'incarico di arbitro è, ovviamente, un argomento tecnico variegato, ma il quadro generale dovrebbe essere piuttosto ovvio: in questo contesto non stiamo parando di condividere il ruolo del GM, ma piuttosto dei modi in cui il gruppo di gioco può portare responsabilità e dialogo creativo nel modo in cui il GM adempie ai suoi compiti.

Quando i compagni di gioco sono in disaccordo su \textit{ruling} puntuali fatti nel bel mezzo della partita, il consiglio autoritario è che tacciano per amore della fluidità nel progresso del gioco. C'è una certa saggezza in questo consiglio, e vi incoraggio a praticarlo come giocatori; lo faccio io stesso, non faccio partire uno stupido dibattito tipo "ad essere precisi, le balestre non funzionano così" ogni volta che ne ho l'occasione. Ci sono un tempo e un luogo per questioni simili.

Il consiglio egalitario che vorrei dirigere all'arbitro, però, è diverso: quando un giocatore è coinvolto nel gioco al punto da contestare un \textit{ruling}, non dirgli di attaccarsi perché tu sei l'arbitro e sta solo distraendo il tavolo. Piuttosto, fai una piccola pausa nell'azione. Chiedigli di spiegare il suo ragionamento, perché pensa che dovresti prendere una decisione diversa. Quindi spiega il tuo ragionamento. E poi \textbf{lasciagli prendere la decisione}, ora che abbiamo fatto lo sforzo di mettere in piedi questo tribunale improvvisato. Cosa ti cambia come si risolve la situazione attuale, finché la decisione è informata ed autentica?

Non voglio fingere che prevenuti interessi umani non contaminino questo tipo di confronti in gioco. Un giocatore autenticamente coinvolto emotivamente potrebbe voler vedere una decisione in proprio favore; non è che voglia barare, è che si concentra sul momento e la propria difesa, non sul quadro generale. La priorità del ruolo dell'arbitro, però, è di favorire l'interazione autentica, secondariamente favorire il principio di diritto e solo al terzo posto assicurarsi che l'\textit{empasse} regolistico corrente si risolva in maniera corretta. Da questa ragione discende il principio dello scambio giudiziale negoziato:

\textbf{Preferisci sempre il riconoscimento di un principio generale a una decisione puntuale; preferisci sempre il coinvolgimento a una regola.}

Alcuni esempi:

\paragraph{Cedi il punto, ottieni conoscenza} Il gruppo cade in una trappola. Un giocatore sostiene che il suo personaggio non fosse lì, ma ammette di non averlo detto ad alta voce perché non era familiare con le procedure del tavolo.

L'arbitro egalitario accetta facilmente l'affermazione in buona fede, ma discute la procedure e si assicura che, in futuro, il gruppo sappia come indicare quando un personaggio devia dall'ordine di marcia. Ha ceduto su una decisione puntuale, affermando il principio generale.

\paragraph{Cedi sulla regola, ottieni partecipazione} Un giocatore è un grande fan dell'arco lungo inglese (una delle opinioni popolari sul medioevo) e ritiene che le regole della campagna non ne rappresentino adeguatamente la bellezza. Per di più, questa insoddisfazione emerge di tanto in tanto in momenti diversi, quindi non è semplicemente un'idea estemporanea.

L'arbitro egalitario riconosce che quale che sia la sua opinione sulla simulazione degli archi, e che egli sia o meno meglio informato del giocatore, non ha tanta importanza quanto il possesso sociale da parte del giocatore della pratica condivisa di giocare assieme; poiché per l'arbitro questa è una piccolezza tra le molte coinvolte in una campagna, mentre per il giocatore l'arco lungo potrebbe essere il punto più interessante dell'intero esercizio, considerato quanto si impegni per sollevare l'argomento! Cedendo su questa piccola cosa, l'arbitro ottiene molto in capitale creativo.

Il principio dello scambio giudiziale negoziato abbandona l'idea dell'arbitro come la bussola senza errore per la direzione creativa della campagna. In pratica svolgerai ancora questo ruolo nove volte su dieci, perché è come si svolge il lavoro creativo, ma questa realtà di base non deve essere affermata tramite pignoleria su discussioni da cortile scolastico quali "naaah, avevo detto chiaramente che ero qua e non là". Cedi, graziosamente, sui particolari, per affermare i valori fondamentali a cui dovreste puntare tutti assieme. Insegna che dire la propria paga, che la risposta dell'arbitro è coinvolta e grata invece che sminuente. Sul lungo termine, ti ripagherà sotto forma di giocatori più coinvolti, intelligenti e ricchi di opinioni.

\usubsection{Prendere sul serio i PX}

Quello che segue suonerà mistico a coloro che sono completamente fuori dal giro e impossibile a coloro che lo sono meno, ma riporerò solamente ciò che so: i punti esperienza nello stile \wargam{e} sono sacri, esistono inter-soggettivamente e sono importanti nello stesso modo in cui lo sono l'handicap nel golf o la classifica in un gioco cabinato. Realizzarlo e trattarli di conseguenza rende il gioco migliore, perché influenza il tuo personale orientamento creativo nel gioco. Essere cinico o, ancora peggio, intenzionalmente non-legittimo, a riguardo riduce la sportività.

Rispettare i punti esperienza significa riconoscere che sono una misura basata su obbiettivi dei risultati nel gioco effettivo: è impossibile ottenerli o perderli per motivi simulativi o ragioni esterne al gioco, perché tracciano l'effettiva performance in gioco del giocatore. I punti sono prezioni perché \dnd{} in stile \wargam{e} è un gioco difficile; il tuo punteggio migliore è un risultato e il tuo punteggio corrente sotto forma dei personaggi ancora vivi è prezioso. Chiunque abbia fatto lo sforzo di portare un personaggio al secondo livello di una campagna legittima sa cosa voglio dire, per il sudore della sua fronte.

L'idea non autentica che i PX siano arbitrari, che l'arbitro "semplicemente decida" chi avanza e quanto velocemente è tossica per il corretto funzionamento del gioco. Personaggi che iniziano direttamente con PX che non si sono guadagnati sono un anatema. Non chiedo a nessuno di credere a qualcosa che non sia vero, ma se non è vero allora il mio consiglio è di renderlo tale: basta essere arbitrari, basta essere accomodanti, basta coprire l'effettiva performance con punti pietà. Solo così si fa sentire che il gioco valga qualcosa emotivamente, fissando il punteggio così che non sia teatro e sia, invece, tanto reale quanto qualsiasi sport. Basta avanzare di livello solo perché il tempo passa.

Siccome i PX sono basati su degli obiettivi (con qualche attività tradizionale estremamente minore come eccezione), si guadagnano in un modo piuttosto difficile che rende misure dirette come "quanto dovresti elargire per sessione" impossibili. La risposta è che dovresti dare zero punti per sessione perché nessuno dovrebbe ricevere anche solo un singolo punto solo per essersi presentato e seduto. Salva una principessa o qualcosa del genere sei vuoi dei PX.

È probabilmente piuttosto ovvio perché, anche se la penso così, do una certa importanza all'avere una serie storica stabile; c'è valore nel poter comparare i punti da sessione a sessione, da campagna a campagna, da GM a GM. Penso anche che le realtà del gioco supportino questo ideale in maniera sorprendentemente stabile. Scambi con altri viaggiatori sembrano confermare che esiste un certo senso di quanto rappresentino, diciamo, mille PX. C'è una certa base per il concetto di legittimità del punteggio, tanto quanto c'è nel golf, nel pattinaggio artistico o nelle classifiche scacchistiche. È piuttosto conveniente abbandonare gli standard storici e fare di testa tua, ma è anche estremamente possibile appoggiarsi a una scala comune, se desideri farlo.

\usubsection{Cos'è la sportività}

La sportività è un concetto piuttosto antiquato e, sebbene veda spesso gente evocarne il nome nel mondo sportivo, raramente vedo qualcuno spiegarne la natura in termini tecnici. Scegliendo lo stile del \wargam{ing}, scegliamo di godere di quest'antica virtù e, dunque, occorre che ne dia una rapida introduzione.

I giochi sono \textbf{spazi rituali} in cui i partecipanti entrano sotto un'assunzione di \textbf{buona fede}. La sportività è una forma di onore che si mantiene con la propria condotta nello spazio rituale. È ciò che permette che l'interazione in gioco sia nettamente positiva; dove altrimenti regnerebbe un'oppressione sociale, il gioco permette apprendimento e crescita personale. Ma questo richiede un comportamento corretto da parte dei partecipanti.

\paragraph{Rispetta gli altri giocatori}: la sportività dà per scontato che non ci sia amicizia tra compagni di tavolo. Dimostrare il contrario è un'infrazione dell'onore non meno grave di violare l'ospitalità, poiché sei entrato nello spazio di gioco sotto mentite spoglie; è qualcosa che fai ai tuoi nemici, un tradimento. Non giochi coi tuoi nemici; giocare è, infatti, un gesto di fiducia e volontà di amicizia. Non trattare i tuoi compagni di gioco come nemici significa garantire loro rispetto: dar loro spazio, ascoltarli e permettere che abbiano la propria agentività.

\paragraph{Apprezza la fiducia}: quando altri sono disposti a giocare con te, tieni il loro cuore tra le tue mani, anche se in un modo limitato. Essere tracotante nella vittoria o amaro nella sconfitta è una violazione della fiducia che ti è stata data quando hanno accettato di giocare con te. Controlla il tuo comportamento, gli altri giocatori non hanno accettato di pulire dopo il tuo passaggio emotivo.

\paragraph{Si vince assieme}: le partite non sono attività a somma zero, dove la vittoria di uno è la sconfitta di un altro. Interagite assieme in gioco per imparare abilità e costruire capitale sociale. Agire diversamente significa diventare uno strumento di sfruttamento o un miope razziatore. Sentirsi meglio a spese degli altri è un comportamento miope e sociopatico. Infangare le tue cerchie sociali chiedendo ai tuoi amici di supportare un comportamento disonorevole è persino peggio.

Lo spazio sociale appropriato per giocare permette ai giocatori di partecipare in maniere che sarebbero pericolose e disinibite negli spazi sociali ordinari. Vi è permesso di perdere senza perdere la faccia, vincere senza deprivare i tuoi pari e imparare senza sottomissione. Questo però dipende dal fatto che i giocatori mantengano le attutidini sociali chiamate "sportività". Apprendetela, è un'abilità fondamentale dell'adulto in una società complessa. Solo i bulli possono andare avanti senza imparare a giocare.

Tutto questo può convincere che i problemi di sportività siano un problema ricorrente in \dnd{}, ma non è davvero così; di solito non se ne parla affatto e per la gran parte delle volte ci organizziamo efficacemente tra noi senza alcun problema. Avere un leader socialmente esperto nel gruppo aiuta, come aiuta avere un gruppo di gioco genericamente adulto, come ci si può aspettare. Ci sono un paio di deboli scenari "da non fare" che sembrano emergere nella cultura di gioco, quindi menzionarli come problemi di sportività da tenere d'occhio potrebbe essere utile per percepire il quadro generale. Impara a non fare queste cose tu stesso, e come interrompere questi pattern comportamentali quando li vedi negli altri:

\paragraph{Amarezza nella sconfitta}: \dnd{} si gioca avendo come avversari scenari arbitrari e la sorte, piuttosto che altri giocatori. Talvolta i giocatori danno la colpa dei loro rovesci all'arbitro, che è francamente non OK quando questo compagno di tavolo ha messo così tanto lavoro nella preparazione del gioco. Cercare di farlo sentire in colpa per qualcosa di ordinario (si perde un sacco in \dnd{}!) non è sportivo. Perché dovrebbe voler arbitrare una partita per te, se la ricompensa è che tu gli scarichi addosso la tua frustrazione per aver fatto del suo meglio?

\paragraph{Sminuire gli altri giocatori}: il gruppo di gioco di \dnd{} è una squadra, con dinamiche da squadra. A volte i giocatori reagiscono a questo tipo di ambiente bullizzando gli altri; per quel che ne so, questa è la programmazione sociale di base del nostro cervello da scimmie, una scorciatoia verso l'unificazione del gruppo tramite la determinazione di un gruppo "in" a spese di un capro espiatorio.

Riconosci questo comportamento (non è una cosa triviale!) e fermalo consicamente. Fare diversamente significa tradire il codice della sportività. Il codice non riconosce alcun tipo di stato da pseudo-partecipante, per cui permetti a una persona di sedere e dividere la pizza con te, ma puoi comunque opprimerlo. Nemmeno se avere un capro espiatorio nella propria cerchia sociale è una vera delizia per il cervello scimmiesco. Nemmeno se quel giocatore è convenientemente marchiato come capro espiatorio da razza, genere, religione, BMI, autismo relativo o qualsiasi altro stato sociale. Devi essere meglio di così per raggiungere la vera sportività.

\paragraph{Insistere di aver ragione}: un giocatore che pensa che avere ragione sia più importante di essere rispettoso non capisce la sportività. \dnd{} richiede un sacco di discussioni, quindi trattare qualsiasi discussione come una minaccia sociale o semplicemente essere ossessivamente concentrati su un solo aspetto porrà rapidamente fine al gioco. Impara a lasciar correre, impara a sbagliarti e impara a modulare il disaccordo; il compromesso non è disonorevole è un segno di rispetto per l'altra parte come partecipante anche quando non sei d'accordo. La verità vincerà comunque sul lungo termine, ma solo se rimani conscio del contesto sociale e accetti di non essere d'accordo sul breve termine.

\paragraph{Rifiutare il supporto emotivo}: questo è difficile, perché le nostre attitudini sociali variano molto tra culture e sottoculture. Non puoi dare regole ad altre persone sul loro coinvolgimento emotivo e farlo provocherà solo altra rabbia. Sei un giudice migliore di me di cosa è sportività costruttiva nell'aiutare i tuoi compagni di gioco a gestire le proprie emozioni di frustrazione, dubbio e rabbia generate dal praticare sport.

Il principio generale è di non abbandonare la gente, non incolparla per le proprie emozioni (interagire seriamente col gioco è una cosa buona!) e tenere aperte le linee di comunicazione. Nella mia esperienza i giocatori preferiscono salvare la faccia tagliando qualsiasi comunicazione, lasciando ogni giocatore a gestire le proprie emozioni da solo e fingendo attivamente di non vedere nulla per eliminare la vergogna. Se puoi fare più di così, si vedranno miglioramenti nella sportività della partita.

\paragraph{Affermare una falsa coscienza}: i gruppi sociali mantengono verità auto-evidenti per i propri membri e un gruppo di gioco non è diverso; tramite interazioni chiuse, formate una camera dell'eco che rinforza le idee condivise. Per evitare di creare e mantenere uno spazio intellettuale isolato, un gruppo dovrebbe cercare di incoraggiare il pensiero critico, la flessibilità intellettuale e una varietà di punti di vista. Fallire in questo significa rischiare di formare una cricca sociale che perde il contatto con la realtà esterna.

Sebbene l'ethos della sportività (o, in questo caso, credo della \wargam{ità}) sia di per sé non-politico, la virtà dell'eccellenza che forma la causa fondamentale dell'attività richiede niente meno che una dedizione alla realtà. Un vero \wargam{e} non può svolgersi a meno che il vostro gruppo di lavoro non sia capace di agire come centro di compensazione per le idee. In questo senso, il tavolo da gioco è come un saloon culturale o uno spazio accademico; non ci si trova per predicare una verità stabilità, ma piuttosto per scoprire conclusioni.

Questa roba da codice sociale è tutto meno che facile, gli adulti falliscono regolarmente nella vita vera quando si tratta di far rispettare questo tipo di standard di società. Essere consci dei problemi e fare del proprio meglio è abbastanza nel senso che parte della ragione per cui giochiamo in generale è per imparare facendo. La sportività è fondamentalmente una serie di virtù sociali che l'attività sportiva cerca di praticare e lo stesso vale per il \wargam{ing}.

\usubsection{La cultura in \dnd{}}

Il \textit{gaming}\translatornote{Mi trovo nuovamente a non avere una traduzione per il termine. Questa volta, però, Eero spiega l'inesistenza di un termine corretto più avanti nella stessa sezione.} come sottocultura è molto legato ai dettagli storici della società di classe americana. Il \textit{gaming} è, specificamente, tanto di classe inferiore (cioè cultura nuova dell'epoca moderna, non educazione ricevuta) e esotico (cioè non parte della cultura comunemente compresa), un esempio primario del concetto americano di "geekiness"\translatornote{Si potrebbe tradurre come "stramberia", ma in questo contesto è un concetto legato alla cultura americana che non ho ritenuto opportuno tradurre.}. La \textit{geekiness} diventa prominente all'inizio del ventesimo secolo, quando la cultura popolare in generale emerge, riflettendo la lotta pre-esistente per definire la cultura normativa dell'alta società. Il \textit{gaming} scivola naturalmente nella scatola \textit{geek} verso la metà del secolo, assieme ai fumetti e a altre perversioni da annusa-mutande.

Il fatto che la storia soprastante sia interamente esclusiva degli Stati Uniti mostra quanto sia complesso il mondo. Culture diverse, con le loro diverse traiettorie attraverso il 20° secolo, possono, al massimo, copiare le attitudini americane in un maniera superficiale. Scimmia vede, scimmia fa. Qui, nella Finlandia post-classicista (accettiamo questa semplificazione per amore del quadro generale, va bene?) la cultura del \textit{gaming}, quando è arrivata nei primi anni '80, si è posizionata come uno di quegli hobby letterari che piacciono studenti universitari e nerd eruditi (il vernacolo finlandese non ha una parola per "geek"). Qualcosa a cavallo tra teatro e letteratura, ma abbastanza giovane, dinamica e punk perché gli adulti siano felici di lasciarti andare avanti per primo.

Io stesso giocavo da una decina d'anni quando ho scoperto su Internet che i giocatori di ruolo americani si identificano generalmente come classe operaia (vale a dire, non abbastanza bravi) e considerano il proprio hobby come infantile e leggermente vergognoso. "Fingere di essere un elfo", come diceva memorabilmente una campagna marketing di \dnd{} degli anni '00\translatornote{Si tratta di una pubblicità stampata per la terza edizione di \dnd{}, che prendeva in giro World of Warcraft. Se è mai arrivata in Italia, non ne ho trovato traccia, ma nemmeno l'originale inglese ha lasciato molti segni. Per i curiosi, il testo completo recitava "If you're going to sit in your basement pretending to be an elf, you should at least have some friends over to help" (se intendi startene seduto in cantina a far finta di essere un elfo, dovresti almeno invitare qualche amico ad aiutarti).}.

L'implicazione, chiara da una prospettiva esterna, è che la cultura del gioco, tradizionalmente e abitualmente, si imbarazza da sola per conformarsi al concetto di \textit{geekiness}. Non lo dico solo riguardo i GDR in comparazione con il resto della "merda \textit{geek}", ma piuttosto per tutta: svegliatevi, siete voi la civiltà occidentale adesso, nel bene e nel male. La supposta alta cultura è morta nel campo della Somme, o così si dice; quel che resta non sono altro che ossa sterili che innalzate sopra di voi. È ora che raccogliate la lanterna della saggezza voi stessi, se qualcuno lo fa.

Nel caso vi stiate chiedendo che cosa tutto questo abbia a che fare con la virtù dell'erudizione in \dnd{}, penso che \dnd{} sia esattamente il tipo di progetto culturale che è stato più colpito dall'istituzione della \textit{geekiness}: per farla breve, come puoi essere ambizioso nella tua erudizione se credi che l'hobby sia intellettualmente triviale?

Ho detto molte cose critiche nei confronti di Gary Gygax in questo libro, \textit{en passant} (non il nostro argomento), ma spezzerò una lancia in suo favore: lui vedeva, credeva e diceva in maniera costante che il gioco di ruolo avesse valore e che valessa la pena che adulti ambiziosi lo prendessero come un serio esercizio intellettuale. Certamente, si è venduto come un vecchio menestrello, ma il \dnd{} che praticava personalmente negli anni '70 sembra essere stato una seria attività di \wargam{ing}, non un gioco per bambini.

\usubsection{Il programma scolastico in breve}

Un altro grande tema in un libro pieno di grandi temi, ma posso quantomeno tracciare i segni principali: questo è quello che vedo quando guardo \dnd{} e ciò che il gioco mi richiede di studiare e comprendere attraverso la sua pedagogia da \wargam{e}, orientata all'attività:

\paragraph{Scienza militare teorica}: Il gioco è inerentemente distaccato dallo studio di un qualsiasi contesto storico specifico. Ciò che resta, quindi, sono i patter strutturali che stanno al di sotto dei conflitti in generale.

\paragraph{Società premoderna}: Visibilmente radicata nella comprensione del pop anni '70 di cosa fosse "medievale", quello di cui il gioco parla davvero è la nostra comprensione dell'umanità premoderna in generale.

\paragraph{Miti, leggende, religione}: Come funziona la magia? Quali sono le leggi delle fate? Il gioco viaggia attraverso densi strati di conoscenza culturale in maniere complesse che richiedono maestria.

\paragraph{Letteratura d'avventura}: semplicemente, un generale \textit{tour de force} letterario che scava nel modo in cui scriviamo e comprendiamo la letteratura e riassembla quel materiale in scenari di gioco.

Questo non è niente, e non so come altro definirlo se non un interesse generale con la virtù dell'"erudizione". Giocare bene a \dnd{} è studiare, giocare e apprendere.